\begin{lemma}
Assume the following main-independence hypothesis: for every real
\(\varepsilon>0\), all sufficiently large \(n\) admit a simple graph \(G\) on
\(\mathrm{Fin}(n)\) with \(G.\mathrm{CliqueFree}(3)\) and
\[
\noderef{IndependenceNumberLess}
\left(G,(1+\varepsilon)\sqrt{(n:\mathbb R)\log(n:\mathbb R)}\right).
\]
Then, for every real \(\eta>0\), there is an integer \(k_0\) such that for
every integer \(k\ge k_0\) there is an integer \(N\) satisfying
\[
\noderef{RamseyWitness}(N,k)
\qquad\text{and}\qquad
\left(\frac12-\eta\right)\noderef{RamseyScale}(k)\le (N:\mathbb R).
\]
\end{lemma}
\begin{proof}
Fix \(\eta>0\).  We prove the displayed eventual statement by cases.

First suppose \(\eta\ge1/2\).  Choose \(k_0=2\).  Let \(k\ge k_0\).  We take
\(N=0\), and let \(G\) be the unique simple graph on \(\mathrm{Fin}(0)\).  This
graph is triangle-free: a triangle would require three vertices of
\(\mathrm{Fin}(0)\), and there are none.

We now check the independent-set part of
\(\noderef{RamseyWitness}(0,k)\) directly from the definition of
\(\noderef{RamseyWitness}\).  Let \(I\) be any finite subset of
\(\mathrm{Fin}(0)\).  Since \(\mathrm{Fin}(0)\) is empty, \(I=\varnothing\) and
\(|I|=0\).  Since \(k\ge2\), no such \(I\) has cardinality \(k\).  Hence there
is no independent set of cardinality \(k\) in \(G\), regardless of the
independence predicate.  Together with triangle-freeness, the graph \(G\)
witnesses \(\noderef{RamseyWitness}(0,k)\).

It remains in this case to verify the scale inequality.  Since \(k\ge2\), we
have \(\log k>0\).  By the definition of \(\noderef{RamseyScale}\),
\[
\noderef{RamseyScale}(k)=\frac{k^2}{\log k}\ge0.
\]
The hypothesis \(\eta\ge1/2\) gives \(1/2-\eta\le0\).  Multiplying the
nonnegative number \(\noderef{RamseyScale}(k)\) by this nonpositive factor gives
\[
\left(\frac12-\eta\right)\noderef{RamseyScale}(k)\le0=(N:\mathbb R).
\]
Thus the required conclusion holds when \(\eta\ge1/2\).

Now assume \(0<\eta<1/2\).  Put
\[
c=\frac12-\frac{\eta}{2}.
\]
Then
\[
\frac12-\eta<c<\frac12,
\qquad
\delta:=c-\left(\frac12-\eta\right)=\frac{\eta}{2}>0.
\]
Because \(2c<1\), choose \(\varepsilon>0\) so small that
\[
2c(1+\varepsilon)^2<1. \tag{1}
\]
Apply the assumed main-independence hypothesis to this \(\varepsilon\).  We
obtain an integer \(n_0\) such that every \(n\ge n_0\) admits a triangle-free
graph \(G\) on \(\mathrm{Fin}(n)\) satisfying
\[
\noderef{IndependenceNumberLess}\!\left(
G,\,(1+\varepsilon)\sqrt{(n:\mathbb R)\log(n:\mathbb R)}\right). \tag{2}
\]

We now choose one eventual threshold for \(k\).  Since \(\log k\to\infty\) and
\[
\frac{k^2}{\log k}\to\infty,
\]
there is an integer \(k_0\) such that for every \(k\ge k_0\):
\[
k\ge3,\qquad \log k>0,\qquad \log k>c,
\]
\[
c\,\frac{k^2}{\log k}\ge n_0+1,\qquad
c\,\frac{k^2}{\log k}\ge4,\qquad
\delta\,\frac{k^2}{\log k}\ge1. \tag{3}
\]
The existence of this single \(k_0\) follows by taking the maximum of finitely
many thresholds: the first three conditions are eventual positivity of
\(\log k\), and the last three use the divergence of \(k^2/\log k\).

Fix \(k\ge k_0\), and abbreviate
\[
X=c\,\frac{k^2}{\log k}.
\]
The conditions in (3) imply \(X\ge0\).  Define the integer
\[
N=N(k)=\lfloor X\rfloor. \tag{4}
\]
For a nonnegative real number \(X\), the floor inequalities give
\[
N\le X=c\,\frac{k^2}{\log k}, \tag{5}
\]
\[
N\ge X-1. \tag{6}
\]
The last inequality in (3) says
\[
1\le \delta\,\frac{k^2}{\log k}.
\]
Subtracting this bound for \(1\) from \(X=c\,k^2/\log k\), we get
\[
X-1\ge
c\,\frac{k^2}{\log k}-\delta\,\frac{k^2}{\log k}
= (c-\delta)\frac{k^2}{\log k}
=\left(\frac12-\eta\right)\frac{k^2}{\log k}. \tag{7}
\]
Combining (6) and (7),
\[
N\ge
\left(\frac12-\eta\right)\frac{k^2}{\log k}. \tag{8}
\]
Similarly, the two middle inequalities in (3) give
\[
N\ge X-1\ge n_0
\qquad\text{and}\qquad
N\ge X-1\ge3. \tag{9}
\]
By \(\noderef{RamseyScale}\), (8) is exactly
\[
\left(\frac12-\eta\right)\noderef{RamseyScale}(k)\le (N:\mathbb R).
\tag{10}
\]

It remains to check that the graph supplied by (2) for this value of \(N\) has
no independent set of cardinality \(k\).  Because \(N\ge3\), we have \(N>0\)
and \(\log N>0\).  From (5) and \(\log k>c>0\), we have
\[
N\le c\,\frac{k^2}{\log k}<k^2. \tag{11}
\]
Thus \(0<N\le k^2\).  The logarithm is increasing on positive real numbers, so
\[
\log N\le \log(k^2)=2\log k. \tag{12}
\]
Both \(N\) and \(\log N\) are nonnegative.  Multiplying the two upper bounds
(5) and (12), and using \(\log k>0\), gives
\[
N\log N
\le
\left(c\,\frac{k^2}{\log k}\right)(2\log k)
=2c\,k^2. \tag{13}
\]
Since \(1+\varepsilon>0\), multiplication by \((1+\varepsilon)^2\) preserves
the inequality:
\[
(1+\varepsilon)^2N\log N
\le 2c(1+\varepsilon)^2k^2
<k^2, \tag{14}
\]
where the final strict inequality is (1).  The left side of (14) is
nonnegative, and \(k>0\).  By monotonicity of square root on nonnegative reals,
and by
\[
\sqrt{(1+\varepsilon)^2N\log N}
=(1+\varepsilon)\sqrt{N\log N},
\]
we obtain
\[
(1+\varepsilon)\sqrt{N\log N}<k. \tag{15}
\]

Now apply (2) at \(n=N\), which is allowed by \(N\ge n_0\) from (9).  We get a
triangle-free graph \(G\) on \(\mathrm{Fin}(N)\) such that
\[
\noderef{IndependenceNumberLess}\!\left(
G,\,(1+\varepsilon)\sqrt{N\log N}\right)
\]
holds.  Unpacking \(\noderef{IndependenceNumberLess}\), every independent
finite set \(I\subseteq \mathrm{Fin}(N)\) satisfies
\[
(|I|:\mathbb R)<(1+\varepsilon)\sqrt{N\log N}. \tag{16}
\]
If there were an independent set \(I\) of cardinality \(k\), then (16) and
(15) would give
\[
(k:\mathbb R)=(|I|:\mathbb R)
<(1+\varepsilon)\sqrt{N\log N}<k,
\]
which is impossible.  Therefore \(G\) has no independent set of cardinality
\(k\).  Together with the triangle-free property from (2), this is precisely a
witness for \(\noderef{RamseyWitness}(N,k)\).  The required scale lower bound
is (10).  Since the same threshold \(k_0\) works for every \(k\ge k_0\), the
desired eventual statement follows.
\end{proof}
